\chapter{DPF Literature-to-Debugging Verification Contract}
\label{ch:bf-dpf-debug-crosswalk}

The preceding chapter turned the DPF and transport staircase into a target-status
taxonomy.  The present chapter turns that taxonomy into a verification
contract.  Its purpose is not to validate any implementation by prose.  Its
purpose is to tell the reader which equation, controlled example,
implementation quantity, tolerance, and failure interpretation belongs to each
layer.

The contract is deliberately ordered.  Target and value-gradient checks come
before tuning.  Teacher-map diagnostics come before learned-map diagnostics.
Equation-local failures should be interpreted narrowly: a failed Sinkhorn
marginal test is evidence about a finite solver, not evidence that particle
flows are wrong; a learned-map residual is evidence about the student and its
training distribution, not evidence that the OT teacher is invalid.

\section{Diagnostic order}
\label{sec:bf-dpf-debug-order}

Use this order before changing code or tuning numerical parameters:
\begin{enumerate}[label=(\arabic*)]
  \item name the scalar object being estimated, relaxed, learned, or sampled;
  \item verify that the gradient path corresponds to that same scalar object
    whenever HMC or autodiff claims are made;
  \item test the classical particle-filter recursion and likelihood estimator
    on a controlled model;
  \item test the particle-flow endpoint and homotopy derivative before testing
    PF-PF proposal correction;
  \item test PF-PF weights and log-determinant accounting before interpreting
    resampling or OT behavior;
  \item test categorical, soft, EOT, Sinkhorn, and barycentric objects before
    testing a learned student map;
  \item only after those checks pass, run posterior comparisons and sampler
    diagnostics.
\end{enumerate}
This order prevents a common error: making a numerically smooth path smoother
without first proving that it is the intended path.

\section{Proposal log-probability autodiff topology}
\label{sec:bf-dpf-proposal-logprob-ad-topology}

The most expensive debugging lesson in one same-object implementation tie-out
was not a floating-point tolerance issue and not a seed issue.  It was an
autodiff-topology issue in the adapted proposal log probability.  The forward
numbers can be exactly right while the derivative is wrong if the proposal
sample and the proposal log probability are wired through the same distribution
object in the computation graph.

Write
\[
  f_\theta(x_t\mid x_{t-1})
  \quad\text{and}\quad
  g_\theta(y_t\mid x_t)
\]
for the transition and observation densities.  For an adapted proposal
\[
  q_\theta(x_t\mid x_{t-1},y_t),
\]
the one-step log-weight update corresponding to
\eqref{eq:bf-pf-sis-recursion} is
\begin{equation}
\label{eq:bf-dpf-adapted-proposal-log-weight}
  \log \tilde w_t^{(i)}
  =
  \log w_{t-1}^{(i)}
  + \log f_\theta(x_t^{(i)}\mid x_{t-1}^{(i)})
  + \log g_\theta(y_t\mid x_t^{(i)})
  - \log q_\theta(x_t^{(i)}\mid x_{t-1}^{(i)},y_t).
\end{equation}
The last term is not the log probability of a random-number-generation event.
It is the density of the realized particle value under the declared proposal
law.

In the linear-Gaussian adapted-proposal case used by one representative debug
trace, the optimal proposal satisfies
\begin{equation}
\label{eq:bf-dpf-optimal-proposal-identity}
  q_\theta(x_t\mid x_{t-1},y_t)
  =
  p_\theta(x_t\mid x_{t-1},y_t)
  =
  \frac{
    f_\theta(x_t\mid x_{t-1})g_\theta(y_t\mid x_t)
  }{
    p_\theta(y_t\mid x_{t-1})
  } .
\end{equation}
Substituting \eqref{eq:bf-dpf-optimal-proposal-identity} into the signed
increment in \eqref{eq:bf-dpf-adapted-proposal-log-weight} gives
\begin{align}
\label{eq:bf-dpf-optimal-proposal-cancellation}
 &\log f_\theta(x_t\mid x_{t-1})
  + \log g_\theta(y_t\mid x_t)
  - \log q_\theta(x_t\mid x_{t-1},y_t) \nonumber\\
 &\hspace{1cm}
  =
  \log p_\theta(y_t\mid x_{t-1}).
\end{align}
Thus, when the proposal is the exact adapted proposal, the signed increment is
independent of the realized value $x_t$ after cancellation.  Equivalently,
\[
  \frac{\partial}{\partial x_t}
  \left[
    \log f_\theta(x_t\mid x_{t-1})
    + \log g_\theta(y_t\mid x_t)
    - \log q_\theta(x_t\mid x_{t-1},y_t)
  \right]
  =0 .
\]
This zero derivative is a cancellation of three density derivatives.  It is not
obtained by making the proposal-density derivative vanish.

The correct implementation topology therefore has two conceptual steps:
\begin{enumerate}[label=(\roman*)]
  \item sample the proposed particle from the proposal law;
  \item evaluate the proposal density at the realized particle as an ordinary
    density argument.
\end{enumerate}
In pseudocode:
\begin{verbatim}
sample_dist = make_proposal_dist(previous_state, observation)
proposed_particles = sample_dist.sample(seed)

density_dist = make_proposal_dist(previous_state, observation)
proposal_ll = density_dist.log_prob(proposed_particles)
\end{verbatim}
The second distribution has the same numerical parameters as the first, but it
is a fresh density-evaluation node.  Equivalently, one may call a helper such
as \texttt{optimal\_proposal\_log\_prob(...)} that rebuilds the proposal mean
and covariance before applying \texttt{log\_prob}.  This is the topology used
by FilterFlow's executable contract: \texttt{propose} first returns a state
whose \texttt{particles} field contains the sampled particles, and
\texttt{loglikelihood(proposed\_state, state, observation)} then recomputes
the proposal distribution and evaluates its density at
\texttt{proposed\_state.particles}.

The tempting but wrong topology is
\begin{verbatim}
proposal_dist = make_proposal_dist(previous_state, observation)
proposed_particles = proposal_dist.sample(seed)
proposal_ll = proposal_dist.log_prob(proposed_particles)
\end{verbatim}
This can be value-correct: the scalar value of \texttt{proposal\_ll} may match
the fresh-density value exactly.  But the derivative can be wrong because the
autodiff graph is no longer the graph of the mathematical density function
$x_t\mapsto \log q_\theta(x_t\mid x_{t-1},y_t)$.  The proposed particle is now
also the reparameterized sample produced by the same distribution node.  In the
row-173 trace, that wiring made the BayesFilter VJP
\[
  \frac{\partial}{\partial x_t}\log q_\theta(x_t\mid x_{t-1},y_t)
\]
with unit upstream equal to zero, while FilterFlow's density-at-realized-particle
route gave a maximum absolute VJP of
\[
  26.075772828701865 .
\]
The forward proposal log-probability values still agreed.  The value agreement
therefore hid the bug.

The downstream effect was exactly the cancellation failure predicted by
\eqref{eq:bf-dpf-optimal-proposal-cancellation}.  With the naive wiring,
BayesFilter's unit-upstream VJP of
\[
  \log f_\theta+\log g_\theta-\log q_\theta
\]
with respect to the proposed particles had maximum absolute discrepancy
\[
  26.075772828710043 ,
\]
and the normalized previous-log-weight carry VJP had discrepancy
\[
  7.930673258915529 .
\]
After evaluating the proposal log probability through a freshly recomputed
proposal density, the proposal-log-probability VJP matched FilterFlow exactly
in that probe, the signed-sum discrepancy fell to about
\[
  1.49\times 10^{-11},
\]
and the normalized carry discrepancy fell to about
\[
  1.45\times 10^{-11}.
\]
These numbers are evidence from one concrete comparison trace; the general rule
is the topology rule above.

The audit implication is strict.  A DPF implementation may pass value replay,
seed replay, particle replay, and proposal-log-probability value replay while
still failing the derivative contract.  Every adapted-proposal implementation
that is used for gradient comparison must therefore check the proposal
log-probability VJP as a density evaluation at the realized particle, not only
the scalar log-probability value.

\section{Same-object deterministic tie-out contract}
\label{sec:bf-dpf-bf-filterflow-deterministic-tieout}

A deterministic implementation tie-out is a same-object comparison, not an
oracle test.  The only statement it supports is conditional:

\begin{quote}
Under a frozen deterministic comparison contract, two implementations can be
shown to match on declared density, filter-path, fixed-ancestor, and fixed-
branch AD-gradient surfaces.
\end{quote}

The contract freezes the objects that would otherwise make a particle-filter
comparison ambiguous:
\begin{enumerate}[label=(\roman*)]
  \item the model row and physical parameter vector;
  \item density probes for initial, transition, and observation log densities;
  \item initial particles, observations, and additive transition innovations;
  \item the scalar, namely the sum of predictive log normalizers for the path
    comparisons;
  \item branch schedules: no resampling for one path test, and fixed ancestor
    indices for the branch-replay path test;
  \item the gradient scalar and physical knobs for the fixed-branch AD-gradient
    comparison.
\end{enumerate}
The comparison deliberately avoids random-number-generator equality.  When an
ancestor schedule is needed, it is treated as data.  Both sides receive the same
schedule and replay the same branch.

The illustrative comparison suite contained six rows:
\[
\begin{gathered}
\texttt{lgssm\_2d\_h25\_rich},\quad
\texttt{sv\_1d\_h18\_rich},\quad
\texttt{range\_bearing\_4d\_h20\_rich},\\
\texttt{structural\_ar1\_quadratic\_h16},\quad
\texttt{spatial\_sir\_j3\_rk4},\quad
\texttt{predator\_prey\_rk4}.
\end{gathered}
\]
The closed deterministic comparison reported:
\begin{itemize}
  \item density tie-out: all six rows matched, with maximum absolute delta
    \(0\);
  \item deterministic no-resampling path tie-out: all six rows matched, with
    maximum absolute delta \(0\);
  \item deterministic fixed-ancestor path tie-out: all six rows matched, with
    maximum absolute delta \(0\), and the expected one resampling event was
    replayed on each executed row;
  \item fixed-branch AD-gradient tie-out: all five executable gradient rows
    matched on scalar and AD-gradient values, with maximum scalar delta \(0\)
    and maximum AD-gradient delta \(0\).
\end{itemize}
The spatial SIR row was not promoted to the gradient comparison.  It remains a
gradient-contract block under the frozen manifest, not a failure and not a
success.

Finite differences are diagnostic-only in this contract.  They are useful for
finding gross scalar/gradient wiring mistakes, but they are not stable enough
to be a gate for the comparison.  In one gradient artifact the
range-bearing same-implementation finite-difference diagnostic had maximum
reported discrepancy about
\[
  1.16\times 10^{-4},
\]
while the BayesFilter and FilterFlow AD gradients still matched exactly under
the fixed-branch scalar.  This is recorded as a numerical diagnostic, not as a
promotion criterion and not as a veto.

Two robustness diagnostics sit below this gate:
\begin{itemize}
  \item a seeded fixed-ancestor diagnostic generated three frozen ancestor
    schedules and replayed them across the full comparison suite.  All executed
    rows matched with maximum absolute delta \(0\).  This is branch-robustness
    evidence only; it is not a stochastic-resampling distribution claim;
  \item a controlled one-dimensional LGSSM horizon ladder compared \(T=2\) and
    \(T=4\).  Both scenarios had implementation agreement.  The \(T=4\)
    scenario retained a forward residual veto in an older Sinkhorn residual
    diagnostic, so the ladder is explanatory and does not change the main
    deterministic comparison conclusion.
\end{itemize}

The non-claims are part of the contract.  The BayesFilter/FilterFlow tie-out
does not prove filter correctness, does not prove that either implementation is
mathematically correct, does not validate stochastic resampling distributions,
does not validate gradients through random or discrete ancestor selection, and
does not compare to student implementations, tensor-train or sparse-grid
alternatives, dense quadrature, paper tables, GPU behavior, HMC behavior, DSGE
behavior, scalability, deployment, or production readiness.

\section{Equation-to-test matrix}
\label{sec:bf-dpf-debug-equation-test-matrix}

\begin{center}
\scriptsize
\begin{longtable}{@{}p{0.13\textwidth} p{0.16\textwidth} p{0.17\textwidth} p{0.17\textwidth} p{0.15\textwidth} p{0.14\textwidth}@{}}
\caption{Equation-indexed verification matrix for the DPF block.}
\label{tab:bf-dpf-debug-equation-test-matrix}\\
\toprule
Layer & Equation or claim & Controlled example & Implementation quantity & Pass signal & Narrow failure interpretation \\
\midrule
\endfirsthead
\caption[]{Equation-indexed verification matrix for the DPF block, continued.}\\
\toprule
Layer & Equation or claim & Controlled example & Implementation quantity & Pass signal & Narrow failure interpretation \\
\midrule
\endhead
\bottomrule
\endlastfoot
Filtering recursion &
\eqref{eq:bf-pf-predictive-law}, \eqref{eq:bf-pf-filtering-law} &
linear-Gaussian state-space model with Kalman reference &
predict/update law, normalized weights, filtered mean/covariance &
moments match reference within declared Monte Carlo error &
recursion, normalization, or model-interface bug; not an OT or learned-map claim \\
Bootstrap likelihood estimator &
\eqref{eq:bf-pf-average-weight}, \eqref{eq:bf-pf-bootstrap-likelihood-estimator} &
same linear-Gaussian model; fixed seeds and increasing $N$ &
per-period average weights and log likelihood estimator &
Monte Carlo mean approaches Kalman likelihood; variance decreases with $N$ &
particle proposal or weight bug; not evidence about differentiability \\
Homotopy derivative &
\eqref{eq:bf-pff-log-homotopy}, \eqref{eq:bf-pff-log-homotopy-derivative} &
scalar nonlinear observation with analytic log-likelihood derivative &
finite difference in $\lambda$ versus analytic derivative &
central difference and analytic derivative agree under step refinement &
homotopy algebra or likelihood derivative issue; not yet a PF-PF correction issue \\
EDH endpoint &
\eqref{eq:bf-pff-edh-velocity}, \eqref{eq:bf-pff-kalman-endpoint} &
linear-Gaussian recovery with known posterior mean/covariance &
flow endpoint particles, mean, covariance, and endpoint residual &
endpoint moments match Kalman update within flow and Monte Carlo tolerances &
flow ODE or Gaussian-closure implementation issue; not HMC target validity \\
LEDH local endpoint &
\eqref{eq:bf-pff-local-jacobian}, \eqref{eq:bf-pff-ledh-ode} &
synthetic observation with known local Jacobian variation &
particle-local Jacobians, local precision, and endpoint continuity &
local quantities are finite and stable under small perturbations &
local linearization fragility; not proof that global EDH is wrong \\
PF-PF weight &
\eqref{eq:bf-pfpf-preflo-density}, \eqref{eq:bf-pfpf-postflow-density}, \eqref{eq:bf-pfpf-weight} &
synthetic affine flow with known determinant &
proposal density, target density, corrected log weight &
manual density ratio and implemented weight agree &
proposal-density or correction algebra bug; not a resampling diagnosis \\
Log-determinant ODE &
\eqref{eq:bf-pfpf-jacobian-ode}, \eqref{eq:bf-pfpf-logdet-ode}, \eqref{eq:bf-pfpf-logdet-affine} &
affine flow $x_\lambda=A_\lambda x+b_\lambda$ &
integrated log determinant and finite-difference determinant &
sign and magnitude match under step-size refinement &
Jacobian sign, time direction, or integrator issue; not a likelihood theorem \\
Categorical discontinuity &
\eqref{eq:bf-soft-categorical-choice} &
two-particle resampling with weight crossing &
ancestor index path and realized resampled cloud &
jump occurs as expected; no false pathwise gradient is reported &
discrete resampling is being differentiated incorrectly; not proof soft resampling is valid \\
Soft-resampling bias &
\eqref{eq:bf-soft-rule-eqn}, \eqref{eq:bf-soft-mean-preservation} &
two-particle nonlinear test function &
mean preservation and nonlinear test-function error &
affine tests behave as derived; nonlinear bias is visible when expected &
surrogate-law effect; not evidence of categorical-target preservation \\
EOT/Sinkhorn marginals &
\eqref{eq:bf-eot-primal}, \eqref{eq:bf-eot-factorization}, \eqref{eq:bf-eot-first-row-balance} &
small cost matrix with prescribed marginals &
row residual, column residual, cost, entropy, $\varepsilon$, iteration count &
residuals fall below tolerance and are stable under log-domain implementation &
finite-solver or stabilization issue; not proof of posterior invariance \\
Barycentric projection &
\eqref{eq:bf-eot-barycentric} &
small coupling with known column masses &
projected equal-weight support cloud and dimensions &
output has declared shape and column-mass normalization &
projection or shape bug; not a coupling-optimality result \\
Learned-map residual &
\eqref{eq:bf-learned-ot-student-map}, \eqref{eq:bf-learned-ot-residual}, \eqref{eq:bf-learned-ot-equivariance} &
permutation-equivariance test and held-out teacher clouds &
teacher/student residuals, permutation residual, training-distribution tags &
residuals stay below threshold on the declared envelope &
student or distribution-shift issue; not a teacher-map refutation \\
HMC value-gradient contract &
\eqref{eq:bf-dpf-hmc-contract} &
finite-difference check on a fixed scalar target and fixed random seed policy &
scalar value, autodiff gradient, finite-difference gradient, accept/reject value &
all paths refer to the same scalar and gradients agree within tolerance &
target-interface bug; do not tune HMC before resolving it \\
\end{longtable}
\end{center}

\section{Minimal controlled examples}
\label{sec:bf-dpf-debug-controlled-examples}

\begin{longtable}{@{}p{0.18\textwidth} p{0.24\textwidth} p{0.26\textwidth} p{0.22\textwidth}@{}}
\caption{Minimal controlled examples for DPF verification.}
\label{tab:bf-dpf-debug-controlled-examples}\\
\toprule
Example & Required inputs & Expected output & Does not prove \\
\midrule
\endfirsthead
\toprule
Example & Required inputs & Expected output & Does not prove \\
\midrule
\endhead
Linear-Gaussian recovery &
Known transition, observation, covariance matrices, Kalman likelihood, and
Kalman filtering moments. &
Particle recursion, EDH endpoint, and likelihood estimates approach the Kalman
reference under increasing particles or decreasing flow step size. &
Nonlinear DSGE validity or learned-map generalization. \\
Synthetic affine flow &
Invertible affine map, analytic determinant, and known proposal density. &
PF-PF density ratio, weight, and log determinant match hand calculation. &
Correctness for nonlinear flow fields. \\
Scalar nonlinear observation &
One-dimensional observation function with analytic derivatives and controlled
curvature. &
Homotopy derivative and local-linearization diagnostics show expected error as
curvature changes. &
High-dimensional stability. \\
Two-particle soft-resampling bias &
Two particles, weights crossing $1/2$, and nonlinear test function. &
Categorical path is discontinuous; soft path has the predicted surrogate bias. &
Acceptability of soft resampling as an HMC target. \\
Small Sinkhorn problem &
Positive cost matrix, weights, equal target marginal, fixed $\varepsilon$, and
known residual tolerance. &
Row/column residuals and stabilized objective behave consistently with the
solver budget. &
Posterior equivalence to categorical resampling. \\
Permutation-equivariance learned-map test &
Teacher clouds, learned map, and permutation matrices. &
Permuting input rows permutes output rows, and scalar summaries remain
invariant when required. &
Small teacher/student error or target preservation. \\
Finite-difference HMC gradient test &
Fixed scalar target, fixed random seed policy, finite-difference step ladder,
and autodiff gradient. &
Autodiff and finite-difference gradients agree for the scalar used by
accept/reject. &
Convergence of the Markov chain or correctness for another scalar target. \\
\bottomrule
\end{longtable}

\section{Diagnostic specification}
\label{sec:bf-dpf-debug-diagnostic-spec}

Every diagnostic result should be recorded with the following fields:
\begin{itemize}
  \item \textbf{Equation or claim}: label from the chapter, not only a method
    name.
  \item \textbf{Inputs}: model, particles, weights, random seed policy,
    $\varepsilon$, iteration budget, architecture, and scalar target as
    applicable.
  \item \textbf{Expected output}: scalar, vector, matrix, empirical measure,
    residual, or posterior summary.
  \item \textbf{Tolerance}: absolute and relative thresholds, Monte Carlo
    standard-error logic, or convergence slope.
  \item \textbf{Failure interpretation}: the narrow mathematical layer that the
    failure implicates.
  \item \textbf{Non-implication}: what the failure does not prove.
\end{itemize}

\begin{longtable}{@{}p{0.15\textwidth} p{0.21\textwidth} p{0.21\textwidth} p{0.19\textwidth} p{0.16\textwidth}@{}}
\caption{Diagnostic contract by layer.}
\label{tab:bf-dpf-debug-diagnostic-contract}\\
\toprule
Diagnostic & Inputs & Tolerance rule & Failure interpretation & Non-implication \\
\midrule
\endfirsthead
\toprule
Diagnostic & Inputs & Tolerance rule & Failure interpretation & Non-implication \\
\midrule
\endhead
Filtering recursion parity &
Kalman reference, fixed particles, fixed model. &
Monte Carlo error decreases with $N$; deterministic quantities match reference
when applicable. &
State-space interface, normalization, or weight update bug. &
Does not reject OT, learned OT, or HMC. \\
Flow endpoint parity &
Linear-Gaussian recovery, flow step ladder. &
Endpoint error decreases as flow integration is refined. &
Flow ODE, closure, or integration issue. &
Does not prove PF-PF correction is wrong. \\
PF-PF algebra parity &
Affine flow, analytic determinant, explicit proposal density. &
Log-weight difference below deterministic tolerance. &
Density-ratio, Jacobian, or sign error. &
Does not diagnose resampling. \\
Soft-resampling bias check &
Two-particle nonlinear test. &
Bias sign and magnitude match derived expression within arithmetic tolerance. &
Surrogate target is being confused with categorical target. &
Does not prove soft resampling is unusable. \\
Sinkhorn residual check &
Small cost matrix, fixed $\varepsilon$, solver budget. &
Row/column residuals below declared threshold and improve with budget. &
Finite-solver or stabilization issue. &
Does not prove EOT is a valid posterior target. \\
Learned residual check &
Teacher outputs, held-out clouds, distribution tags. &
Residuals reported by regime; thresholds tied to training envelope. &
Student approximation or extrapolation issue. &
Does not invalidate the teacher. \\
HMC value-gradient check &
Scalar value, autodiff path, finite-difference ladder. &
Gradient errors shrink under step refinement until numerical limits. &
Value-gradient mismatch or nondifferentiable scalar path. &
Does not prove chain convergence. \\
\bottomrule
\end{longtable}

\section{Source and implementation map}
\label{sec:bf-dpf-debug-source-implementation-map}

\begin{center}
\scriptsize
\begin{longtable}{@{}p{0.14\textwidth} p{0.18\textwidth} p{0.20\textwidth} p{0.18\textwidth} p{0.20\textwidth}@{}}
\caption{Source, chapter, equation, and implementation links.}
\label{tab:bf-dpf-debug-source-implementation-map}\\
\toprule
Layer & Chapter location & Equation labels & Source family & Implementation quantity \\
\midrule
\endfirsthead
\toprule
Layer & Chapter location & Equation labels & Source family & Implementation quantity \\
\midrule
\endhead
Classical PF &
Chapter~\ref{ch:bf-particle-filters}, Sections~\ref{sec:bf-pf-ssm-recursion}--\ref{sec:bf-pf-likelihood-estimator} &
\eqref{eq:bf-pf-filtering-law}, \eqref{eq:bf-pf-bootstrap-likelihood-estimator} &
SMC literature spine &
weights, ESS, likelihood estimator, empirical moments \\
Particle flow &
Chapter~\ref{ch:bf-dpf-particle-flow-foundations}, Sections~\ref{sec:bf-pff-homotopy}--\ref{sec:bf-pff-ledh} &
\eqref{eq:bf-pff-log-homotopy-derivative}, \eqref{eq:bf-pff-edh-ode}, \eqref{eq:bf-pff-ledh-ode} &
Daum--Huang, PF-PF flow literature &
flow field, endpoint residual, local Jacobians \\
PF-PF correction &
Chapter~\ref{ch:bf-dpf-pfpf}, Sections~\ref{sec:bf-pfpf-change-of-variables}--\ref{sec:bf-pfpf-logdet} &
\eqref{eq:bf-pfpf-weight}, \eqref{eq:bf-pfpf-logdet-ode} &
proposal correction and change-of-variables spine &
proposal density, target density, log determinant, corrected weight \\
Resampling and OT &
Chapter~\ref{ch:bf-dpf-soft-resampling}; Chapter~\ref{ch:bf-dpf-ot-baseline}; Chapter~\ref{ch:bf-dpf-entropic-ot-sinkhorn} &
\eqref{eq:bf-soft-categorical-choice}, \eqref{eq:bf-soft-mean-preservation}, \eqref{eq:bf-ot-baseline-lp}, \eqref{eq:bf-eot-primal} &
differentiable resampling, OT, Sinkhorn sources &
ancestor path, soft surrogate, coupling residuals, barycentric cloud \\
Learned OT &
Chapter~\ref{ch:bf-dpf-learned-ot} &
\eqref{eq:bf-learned-ot-student-map}, \eqref{eq:bf-learned-ot-residual}, \eqref{eq:bf-learned-ot-equivariance} &
EOT teacher and set-architecture sources &
teacher output, student output, residuals, distribution tags \\
HMC target &
Chapter~\ref{ch:bf-dpf-hmc-target-suitability}, Sections~\ref{sec:bf-dpf-hmc-contract}--\ref{sec:bf-dpf-hmc-banking-gates} &
\eqref{eq:bf-dpf-hmc-contract} &
HMC, pseudo-marginal, surrogate-HMC spine &
scalar value, gradient, accept/reject value, target-status label \\
\bottomrule
\end{longtable}
\end{center}

\section{Scalable OT reuse-risk and validation doctrine}
\label{sec:bf-dpf-debug-scalable-ot-stack}

The scalable-OT survey sharpens the verification contract for the OT block in
three ways.  First, code availability is useful but does not by itself justify
promotion.  Second, exact-online, retained-teacher approximate, and
structure-changing transport families should not share the same pass/fail story.
Third, the validation ladder must test transported particle objects, not only OT
scalars or paper-reported benchmark values.

A practical BayesFilter review should therefore separate three lanes:
\begin{enumerate}[label=(\roman*)]
  \item \textbf{Exact reference lane:} dense or streamed entropic OT/Sinkhorn
    implementations that preserve the same declared teacher object.
  \item \textbf{Retained-teacher approximation lane:} Nystrom, positive-feature,
    low-rank, or similar accelerators judged first by teacher-preservation
    residuals.
  \item \textbf{Structure-changing lane:} sliced, localized, block, or other
    replacement transport rules judged as explicit new resampling
    approximations.
\end{enumerate}
The first question for any candidate is therefore not ``is it faster?'' but
``which of these lanes does it belong to?''  The rest of the diagnostic stack
depends on that classification.

\subsection{Code availability and reuse risk are prefilters, not verdicts}
\label{sec:bf-dpf-debug-code-risk}

Source availability is a practical filter because BayesFilter needs an exposed
transport object, not only a scalar loss.  But a repository with a readable
commit, a visible API, or even a working prototype does not become default-ready
for BayesFilter on that basis alone.  Reuse risk should therefore be recorded as a
narrow implementation note:
\begin{enumerate}[label=(\arabic*)]
  \item whether the code exposes a usable transport object or only a scalar loss,
  \item whether the backend is compatible with the BayesFilter TensorFlow/TFP
    default,
  \item whether the source is complete and inspectable,
  \item whether the code path matches the mathematical object cited in the
    literature,
  \item whether any missing piece would force BayesFilter to invent new semantics
    locally rather than reproduce the cited object.
\end{enumerate}
A failure here does not refute the mathematics.  It only limits reuse value and
raises implementation-surface risk.

\subsection{Validation ladder before promotion}
\label{sec:bf-dpf-debug-validation-ladder}

Before any scalable OT candidate is promoted beyond a literature note or local
prototype, the review should follow a fixed ladder.  The central rule is that
transported-particle parity comes before downstream performance claims.
\begin{enumerate}[label=(\arabic*)]
  \item deterministic dense parity on small fixtures: compare transported
    particles or transported summaries, not only OT objective values;
  \item marginal checks: source and target mass constraints under weighted inputs;
  \item positivity, non-finite, and numerical-stability checks;
  \item representation diagnostics appropriate to the lane: rank, feature
    dimension, sparsity, projection count, block locality, or minibatch design;
  \item LEDH-PFPF-OT value parity on existing BayesFilter fixtures;
  \item gradient smoke tests only after value-side parity is understood;
  \item runtime and memory ladder over particle count and state dimension;
  \item downstream filtering diagnostics;
  \item multi-seed or uncertainty-aware comparisons before ranking stochastic
    candidates;
  \item default-readiness review only after downstream evidence, not from source
    availability or paper benchmarks alone.
\end{enumerate}
This ladder is deliberately asymmetric.  Passing an early rung does not justify a
later claim, but failing an early rung can already block promotion.

\subsection{Lane-specific narrow interpretations}
\label{sec:bf-dpf-debug-lane-interpretation}

The failure interpretation should remain lane-specific.
\begin{itemize}
  \item In the exact reference lane, a failed marginal or barycentric parity check
    is evidence about the finite solver or implementation path.
  \item In the retained-teacher approximation lane, a large teacher-preservation
    residual is evidence about the approximation family or its tested envelope,
    not a refutation of the teacher object itself.
  \item In the structure-changing lane, disagreement with the dense teacher is not
    automatically failure; the question is whether the replacement criterion and
    downstream validation have been declared and met.
\end{itemize}
This prevents a common review error: treating every non-dense transport as if it
were trying to prove exact parity with the same teacher object.

\section{Promotion thresholds}
\label{sec:bf-dpf-debug-promotion-thresholds}

\begin{longtable}{@{}p{0.19\textwidth} p{0.30\textwidth} p{0.30\textwidth} p{0.13\textwidth}@{}}
\caption{Claim thresholds for DPF implementation interpretations.}
\label{tab:bf-dpf-debug-promotion-thresholds}\\
\toprule
Claim level & Required passing evidence & Disallowed interpretation & Gate \\
\midrule
\endfirsthead
\toprule
Claim level & Required passing evidence & Disallowed interpretation & Gate \\
\midrule
\endhead
Exploratory implementation &
Layer-local smoke tests pass on minimal controlled examples; failures are
recorded with equation labels. &
Treating smoothness or runtime as target correctness. &
May run more experiments. \\
Credible research implementation &
Classical PF, flow, PF-PF, resampling/OT, learned-map, and HMC value-gradient
checks pass on named model classes with sensitivity reports. &
Generalizing beyond tested $N$, $d_x$, $\varepsilon$, model class, or target
status. &
May make bounded research claims. \\
Broader external claim &
Research checks plus posterior comparisons, stress regimes, reproducible
demonstrations, and documented failure policies. &
Claiming deployment suitability from a chapter or benchmark alone. &
Requires separate domain review. \\
Operational deployment claim &
Independent review, change control, monitoring, rollback, data/model
governance, release stability, and ongoing diagnostics. &
Assuming chapter-level mathematics alone is sufficient approval. &
Requires a separate operational process. \\
\bottomrule
\end{longtable}

\section{Boundary of this contract}
\label{sec:bf-dpf-debug-boundary}

This chapter does not validate any DPF implementation.  It defines the
minimum diagnostic evidence needed to interpret failures and to avoid promoting
a method beyond its evidence.  The word ``validate'' is therefore reserved for a
specific equation-local check or for a separately administered external review;
it is not used here as a blanket claim about DPF, learned OT, or HMC.

\FloatBarrier
